% GENERATED FILE. Edit canonical lessons, then run npm run generate:books.
% canonical-source-hash: fnv1a64:f3a1aabc145d9327
% canonical-lessons: MR-C78-dhavne, MR-C78-khelne, MR-C78-vikne, MR-C78-ughadne, MR-C78-shijavne

\chapter{To run, To play, To sell, To open, To cook}
\label{ch:mr-to-run-to-play-to-sell-to-open-to-cook}
\hlchaptermodality{78}

\begin{chapteropening}
\textbf{By the end of this chapter:} \emph{I can say to run, to play, to sell, to open and to cook.}

Complete the last lesson of chapter 78: i can say to run, to play, to sell, to open and to cook.
\end{chapteropening}

\section[dhāvṇe]{\mr{धावणे} (dhāvṇe) — to run}

\label{lesson:MR-C78-dhavne}



\begin{quote}
Before the new one: say the Marathi for to hear, then the Marathi for to get up.
\end{quote}

\subsection*{You'll want to know: \mr{धावणे}}
\textbf{\mr{धावणे}} — \emph{dhāvṇe} — \textquotedblleft{}to run\textquotedblright{}.

\begin{grammarlens}[title={What you've built}]
One more word for this chapter.
\end{grammarlens}

\subsection*{Your turn}
\begin{itemize}
\raggedright
  \item \emph{Say it:} \emph{dhāvṇe}
  \item \emph{Say it:} \emph{dhāvṇe}, once more
  \item \emph{Say it:} say \emph{dhāvṇe}
  \item \emph{Recall:} say the Marathi for to hear, then the Marathi for to get up, then say \emph{dhāvṇe} again
\end{itemize}

\begin{tcolorbox}[breakable,colback=teal!4,colframe=teal!35!black,title={Before you move on}]
What does \mr{धावणे} mean? (\textquotedblleft{}To run\textquotedblright{}.) Say it once more.
\end{tcolorbox}
\section[kheḷṇe]{\mr{खेळणे} (kheḷṇe) — to play}

\label{lesson:MR-C78-khelne}



\begin{quote}
Before the new one: say the Marathi for to get up, then the Marathi for to run.
\end{quote}

\subsection*{You'll want to know: \mr{खेळणे}}
\textbf{\mr{खेळणे}} — \emph{kheḷṇe} — \textquotedblleft{}to play\textquotedblright{}.

\begin{grammarlens}[title={What you've built}]
One more word for this chapter.
\end{grammarlens}

\subsection*{Your turn}
\begin{itemize}
\raggedright
  \item \emph{Say it:} \emph{kheḷṇe}
  \item \emph{Say it:} \emph{kheḷṇe}, once more
  \item \emph{Say it:} say \emph{kheḷṇe}
  \item \emph{Recall:} say the Marathi for to get up, then the Marathi for to run, then say \emph{kheḷṇe} again
\end{itemize}

\begin{tcolorbox}[breakable,colback=teal!4,colframe=teal!35!black,title={Before you move on}]
What does \mr{खेळणे} mean? (\textquotedblleft{}To play\textquotedblright{}.) Say it once more.
\end{tcolorbox}
\section[vikṇe]{\mr{विकणे} (vikṇe) — to sell}

\label{lesson:MR-C78-vikne}



\begin{quote}
Before the new one: say the Marathi for to run, then the Marathi for to play.
\end{quote}

\subsection*{You'll want to know: \mr{विकणे}}
\textbf{\mr{विकणे}} — \emph{vikṇe} — \textquotedblleft{}to sell\textquotedblright{}.

\begin{grammarlens}[title={What you've built}]
One more word for this chapter.
\end{grammarlens}

\subsection*{Your turn}
\begin{itemize}
\raggedright
  \item \emph{Say it:} \emph{vikṇe}
  \item \emph{Say it:} \emph{vikṇe}, once more
  \item \emph{Say it:} say \emph{vikṇe}
  \item \emph{Recall:} say the Marathi for to run, then the Marathi for to play, then say \emph{vikṇe} again
\end{itemize}

\begin{tcolorbox}[breakable,colback=teal!4,colframe=teal!35!black,title={Before you move on}]
What does \mr{विकणे} mean? (\textquotedblleft{}To sell\textquotedblright{}.) Say it once more.
\end{tcolorbox}
\section[ughaḍṇe]{\mr{उघडणे} (ughaḍṇe) — to open}

\label{lesson:MR-C78-ughadne}



\begin{quote}
Before the new one: say the Marathi for to play, then the Marathi for to sell.
\end{quote}

\subsection*{You'll want to know: \mr{उघडणे}}
\textbf{\mr{उघडणे}} — \emph{ughaḍṇe} — \textquotedblleft{}to open\textquotedblright{}.

\begin{grammarlens}[title={What you've built}]
One more word for this chapter.
\end{grammarlens}

\subsection*{Your turn}
\begin{itemize}
\raggedright
  \item \emph{Say it:} \emph{ughaḍṇe}
  \item \emph{Say it:} \emph{ughaḍṇe}, once more
  \item \emph{Say it:} say \emph{ughaḍṇe}
  \item \emph{Recall:} say the Marathi for to play, then the Marathi for to sell, then say \emph{ughaḍṇe} again
\end{itemize}

\begin{tcolorbox}[breakable,colback=teal!4,colframe=teal!35!black,title={Before you move on}]
What does \mr{उघडणे} mean? (\textquotedblleft{}To open\textquotedblright{}.) Say it once more.
\end{tcolorbox}
\section[shijavṇe]{\mr{शिजवणे} (shijavṇe) — to cook}

\label{lesson:MR-C78-shijavne}



\begin{quote}
Before the new one: say the Marathi for to sell, then the Marathi for to open.
\end{quote}

\subsection*{You'll want to know: \mr{शिजवणे}}
\textbf{\mr{शिजवणे}} — \emph{shijavṇe} — \textquotedblleft{}to cook\textquotedblright{}.

\begin{grammarlens}[title={What you've built}]
One more word for this chapter.
\end{grammarlens}

\subsection*{Your turn}
\begin{itemize}
\raggedright
  \item \emph{Say it:} \emph{shijavṇe}
  \item \emph{Say it:} \emph{shijavṇe}, once more
  \item \emph{Say it:} say \emph{shijavṇe}
  \item \emph{Recall:} say the Marathi for to sell, then the Marathi for to open, then say \emph{shijavṇe} again
\end{itemize}

\begin{tcolorbox}[breakable,colback=teal!4,colframe=teal!35!black,title={Before you move on}]
What does \mr{शिजवणे} mean? (\textquotedblleft{}To cook\textquotedblright{}.) Say it once more.
\end{tcolorbox}
